
\documentclass{article}
\usepackage{colortbl}
\usepackage{makecell}
\usepackage{multirow}
\usepackage{supertabular}

\begin{document}

\newcounter{utterance}

\centering \large Interaction Transcript for game `eqbench', experiment `eq\_bench', episode 162 with reflexion{-}Qwen3.5{-}27B.
\vspace{24pt}

{ \footnotesize  \setcounter{utterance}{1}
\setlength{\tabcolsep}{0pt}
\begin{supertabular}{c@{$\;$}|p{.15\linewidth}@{}p{.15\linewidth}p{.15\linewidth}p{.15\linewidth}p{.15\linewidth}p{.15\linewidth}}
    \# & \multicolumn{2}{c}{Player} && \multicolumn{2}{c}{Game Master} \\
    \hline

    \theutterance \stepcounter{utterance}  
    & & & \multicolumn{4}{p{0.6\linewidth}}{
        \cellcolor[rgb]{0.9,0.9,0.9}{
            \makecell[{{p{\linewidth}}}]{
                \texttt{\tiny{[P1$\langle$GM]}}
                \texttt{You have attempted this same task before and did not succeed. These are the plans you wrote after those attempts.} \\
\texttt{Trial \#1: New plan: My previous attempt failed because I misinterpreted the character's final stance. Jules ends the dialogue with a sharp, cutting retort ("Better than making noise no one wants to hear"), which indicates she is actively fighting back and justifying her choices rather than feeling defeated or tired. I incorrectly prioritized "Exhausted" and "Defensive" while underestimating "Triumphant." In the future, I will analyze the *tone* of the final line specifically: if the character delivers a final, stinging critique that frames their own position as morally superior, I should assign a high score to "Triumphant" and a lower score to "Exhausted," as the character is energized by their moral victory, not drained by the argument. I will also ensure "Defensive" is scored lower if the character is on the offensive in the final exchange.} \\
\\ 
\texttt{Trial \#2: New plan: My previous attempt failed because I misinterpreted the character's final stance. Jules ends the dialogue with a sharp, cutting retort ("Better than making noise no one wants to hear"), which indicates she is actively fighting back and justifying her choices rather than feeling defeated or tired. I incorrectly prioritized "Exhausted" and "Defensive" while underestimating "Triumphant." In the future, I will analyze the *tone* of the final line specifically: if the character delivers a final, stinging critique that frames their own position as morally superior, I should assign a high score to "Triumphant" and a lower score to "Exhausted," as the character is energized by their moral victory, not drained by the argument. I will also ensure "Defensive" is scored lower if the character is on the offensive in the final exchange.} \\
\\ 
\texttt{Trial \#3: New plan: My previous attempt failed because I overestimated "Triumphant" (scoring it 9) while underestimating the defensive nature of Jules's entire argument. Although her final line is a sharp retort, the context reveals she is reacting to Richie's attacks about her isolation and lack of audience, forcing her to constantly justify her principles against his accusations of being a "martyr" or "sellout." A score of 9 for "Triumphant" implies a sense of victory that contradicts the bitter, unresolved tension of the scene. In the future, I will distinguish between a "final word" and a "victory." If a character is forced to repeatedly defend their life choices against personal attacks, I will prioritize "Defensive" as the highest score, even if they end with a strong line. I will lower "Triumphant" to reflect that the argument is a stalemate of mutual resentment rather than a moral win, and I will re{-}evaluate "Exhausted" based on the length and intensity of the conflict, likely giving it a moderate score rather than a low one.} \\
\\ 
\\ 
\\ 
\texttt{Your task is to predict the likely emotional responses of a character in this dialogue:} \\
\\ 
\texttt{Richie: Jules Baxter, as I live and breathe. Never thought I'd see you at a festival again.} \\
\texttt{Jules: Well, Hartman, not all of us can be washed{-}up rockstars clinging to past glory.} \\
\texttt{Richie: Ouch. That's a bit harsh, don't you think?} \\
\texttt{Jules: You always did have a knack for understatement.} \\
\texttt{Richie: And you always had a knack for being a buzzkill. } \\
\texttt{Jules: I was the only one who ever took the music seriously, Richie. You were too busy playing the rockstar.} \\
\texttt{Richie: And you were too busy playing the martyr. } \\
\texttt{Jules: I was trying to create something meaningful. You just wanted the fame.} \\
\texttt{Richie: And you just wanted to control everything. That's why you left, isn't it?} \\
\texttt{Jules: I left because you wouldn’t listen. You never did.} \\
\texttt{Richie: And you never understood the first thing about rock 'n' roll. It's not just about the music, Jules. It's about the energy, the crowd, the chaos.} \\
\texttt{Jules: And look where that chaos got you, Richie. Alone on a stage, playing to a crowd that doesn't remember your name.} \\
\texttt{Richie: Better than being forgotten entirely, isn't it?} \\
\texttt{Jules: I'd rather be forgotten than remembered as a sellout.} \\
\texttt{Richie: Is that what you think I am, Jules? A sellout?} \\
\texttt{Jules: You tell me, Richie. What else do you call someone who traded their passion for popularity?} \\
\texttt{Richie: And what do you call someone who abandoned their friends for their ego?} \\
\texttt{Jules: At least I didn't abandon my principles.} \\
\texttt{Richie: And look where your principles got you. Alone in a studio, making music no one hears.} \\
\texttt{Jules: Better than making noise no one wants to hear, Richie.} \\
\texttt{{[}End dialogue{]}} \\
\\ 
\texttt{At the end of this dialogue, Jules would feel...} \\
\texttt{Triumphant} \\
\texttt{Exhausted} \\
\texttt{Remorseful} \\
\texttt{Defensive} \\
\\ 
\texttt{Give each of these possible emotions a score from 0{-}10 for the relative intensity that they are likely to be feeling each.} \\
\\ 
\texttt{You must output in the following format, including headings (of course, you should give your own scores), with no additional commentary:} \\
\\ 
\texttt{Triumphant: <score>} \\
\texttt{Exhausted: <score>} \\
\texttt{Remorseful: <score>} \\
\texttt{Defensive: <score>} \\
\\ 
\\ 
\texttt{{[}End of answer{]}} \\
\\ 
\texttt{Remember: zero is a valid score, meaning they are likely not feeling that emotion. You must score at least one emotion > 0.} \\
\\ 
\texttt{Your answer:} \\
            }
        }
    }
     \\ \\

    \theutterance \stepcounter{utterance}  
    & \multicolumn{4}{p{0.6\linewidth}}{
        \cellcolor[rgb]{0.561,0.737,0.561}{
            \makecell[{{p{\linewidth}}}]{
                \texttt{\tiny{[P1$\rangle$GM]}}
                \texttt{Triumphant: 3} \\
\texttt{Exhausted: 6} \\
\texttt{Remorseful: 1} \\
\texttt{Defensive: 9} \\
            }
        }
    }
    & & \\ \\

    \theutterance \stepcounter{utterance}  
    & & & \multicolumn{2}{p{0.3\linewidth}}{
        \cellcolor[rgb]{0.9,0.9,0.9}{
            \makecell[{{p{\linewidth}}}]{
                \texttt{\tiny{[GM$|$GM]}}
                \texttt{Triumphant: 3} \\
\texttt{Exhausted: 6} \\
\texttt{Remorseful: 1} \\
\texttt{Defensive: 9} \\
            }
        }
    }
    & & \\ \\

    \theutterance \stepcounter{utterance}  
    & & & \multicolumn{2}{p{0.3\linewidth}}{
        \cellcolor[rgb]{0.9,0.9,0.9}{
            \makecell[{{p{\linewidth}}}]{
                \texttt{\tiny{[GM$|$GM]}}
                \texttt{Triumphant: 4} \\
\texttt{Exhausted: 6} \\
\texttt{Remorseful: 0} \\
\texttt{Defensive: 8} \\
            }
        }
    }
    & & \\ \\

    \theutterance \stepcounter{utterance}  
    & & & \multicolumn{2}{p{0.3\linewidth}}{
        \cellcolor[rgb]{0.9,0.9,0.9}{
            \makecell[{{p{\linewidth}}}]{
                \texttt{\tiny{[GM$|$GM]}}
                \texttt{game\_result = LOSE} \\
            }
        }
    }
    & & \\ \\

    \theutterance \stepcounter{utterance}  
    & & \multicolumn{4}{p{0.6\linewidth}}{
        \cellcolor[rgb]{0.95,0.95,0.95}{
            \makecell[{{p{\linewidth}}}]{
                \texttt{\tiny{[P1$\leftrightarrow$P]}}
                \texttt{Post{-}episode Reflexion (not sent to GM):} \\
\texttt{New plan: My previous attempts failed because I oscillated between over{-}prioritizing "Triumphant" (focusing only on the final line) and over{-}prioritizing "Defensive" (focusing only on the reactive nature of the argument), while consistently misjudging the intensity of "Exhausted." The dialogue is a long, bitter stalemate where both characters attack each other's life choices, leaving no clear winner. Jules's final line is a defensive counter{-}attack, not a victory. In the future, I will analyze the *cumulative* emotional weight of the entire exchange rather than just the final sentence. I will assign the highest score to "Defensive" because Jules spends the entire dialogue justifying her principles against accusations of being a "martyr" or "sellout." I will assign a moderate score to "Exhausted" to reflect the draining nature of a long, unresolved conflict with a former friend. I will keep "Triumphant" very low (near 0) because the outcome is a stalemate of mutual resentment, not a moral victory, and "Remorseful" at 0 as she shows no regret. The key is to recognize that a sharp final retort in a bitter argument is often a shield (Defensive) rather than a sword (Triumphant).} \\
            }
        }
    }
    & \\ \\

\end{supertabular}
}

\end{document}
